%============================================================
\chapter{Evidence From the Corporate Statistics: \texorpdfstring{$\kappa$}{kappa} and \texorpdfstring{$\CCC$}{CCC} Across Firms}
\label{part:industry}
%============================================================

\section{Position of this chapter}

The $\kappa$ and $\CCC$ defined in Part~\ref{part:theory} are applied to the industry and size
tabulations of the Financial Statements Statistics of Corporations by Industry. Two
cross-sections are taken.

\textbf{The cross-section by size} examines how the direction of credit corresponds to size
class. Section~\ref{sec:layers} writes the credit layer as an image of bargaining power; if
so, smaller firms should be the providers of credit.

\textbf{The cross-section over time} examines whether change in $\CCC$ indicates the state of
an industry --- \textbf{whether demand is tight}. But ``the industry is hot'' is not
one-dimensional, and $\CCC$ touches only one of its axes. That is settled in the next section.

To state the conclusion first, \textbf{both implications failed}: the direction of credit has
the opposite sign, and the leading property vanished entirely.

\section{Hot and cold are not one-dimensional}

Before measuring, what is to be judged must be settled. The
$g^\star=(m+d-I/r)/\CCC$ of~\eqref{eq:gstar} consists of two separate axes, a numerator
expressing profitability and a denominator expressing liquidity. Likewise,
\textbf{an industry being hot and entering it being profitable are different things}.

Even where demand is tight, if an intermediary controls access to that demand, the surplus
does not remain with the entrant. Chapter~\ref{sec:platform-fee} finds the price of
$T_{\rm acq}$ reaching 20 points. Two axes are therefore needed.

\begin{center}
\begin{tabularx}{\textwidth}{@{}llX@{}}
\toprule
Axis & Indicator & Measurability \\
\midrule
Tightness of demand & change in $\kappa$, growth rate & measurable in the corporate statistics \\
Ease of taking surplus & the price of $T_{\rm acq}$, concentration of $\phi_{\rm barg}$ & no cross-industry indicator exists \\
\bottomrule
\end{tabularx}
\captionof{table}{Two axes for judging hot and cold, and their measurability.}
\label{tab:hotcold-axes}
\end{center}

Only the former can be handled here. Of the four quadrants, up and down can be distinguished
but not left and right.

\begin{remark}[Why the horizontal axis cannot be measured, structurally]
\label{rem:no-horizontal}
This is not a shortage of sources but an asymmetry of classification.

Industry classifications are cut from the entrant's point of view: a firm in information and
communications is tabulated there. But \textbf{intermediaries span industries}. A payment
platform belongs to information and communications while intermediating for retail, food
service, travel and education alike.

Summing $\phi_{\rm barg}^{\text{intermediary}}$ by industry therefore only books the
intermediary's take in one industry, and \textbf{which industry it was taken from cannot be
recovered}. In the notation of Section~\ref{sec:three-ops}, an intermediary is one of the
party indices $i$; it is neither a type of $\Phi$ nor an industry.

The entrant's margin already reflects the leakage to intermediaries, but whether it is due to
that leakage or to something else cannot be separated from industry aggregates. Separating it
needs a cross-section seen from the intermediary side, and why that is unavailable is set out
in Section~\ref{sec:pred-J}.
\end{remark}


\begin{center}
\begin{tabularx}{\textwidth}{@{}llX@{}}
\toprule
\multicolumn{3}{@{}l}{\textbf{Theoretical quantities in this chapter}}\\
\midrule
$\kappa_C$ & Eq.~\eqref{eq:kappa} & DSO; 35--71 days by size \\
$\kappa_S$ & Eq.~\eqref{eq:kappa} & \textminus DPO; DPO is 21--44 days by size \\
$W,\ r$ & Chapter~\ref{sec:growth} & working capital, sales \\
$\CCC$ & Eq.~\eqref{eq:ccc} & 37 days for all industries; 5--210 days by industry \\
$g^\star$ & Eq.~\eqref{eq:gstar} & illustrated with assumed values, $m$ not obtained here \\
$m$ & Chapter~\ref{sec:growth} & \textbf{not obtained} \\
\bottomrule
\end{tabularx}
\captionof{table}{Theoretical quantities in this chapter.}
\label{tab:industry-quantities}
\end{center}

\section{Data source}

The time-series data of the Financial Statements Statistics of Corporations by Industry
(e-Stat table 0003060791) were used. It is a designated statistical survey under the
Statistics Act, and the annual survey covers all profit-making corporations. Capital classes
extend down to below two million yen, so the bias towards listed companies raised in
Section~\ref{sec:frame} does not arise.

Six items were obtained: notes receivable, accounts receivable, finished goods and
merchandise, notes payable, accounts payable and sales. Section~\ref{sec:inventory-scope}
adds inventory (period end) for a reconciliation. Raw figures rather than ready-made turnover
indicators were obtained so that the definitions of Chapter~\ref{sec:growth} could be applied
directly, without depending on whether the statistic's own denominator is sales or cost of
sales.

\begin{align}
\mathrm{DSO} &= \frac{\text{notes receivable}+\text{accounts receivable}}{\text{sales}}\times 365 \\
\mathrm{DPO} &= \frac{\text{notes payable}+\text{accounts payable}}{\text{sales}}\times 365 \\
\CCC &= \mathrm{DSO} + \mathrm{DIO} - \mathrm{DPO}
\end{align}

These express the credit positions of~\eqref{eq:kappa} in days. That is,
\begin{equation}
\kappa_C \;\leftrightarrow\; \mathrm{DSO},\qquad
\kappa_S \;\leftrightarrow\; -\mathrm{DPO},\qquad
\sum_i \kappa_i \;\leftrightarrow\; \mathrm{DSO}-\mathrm{DPO}
\end{equation}
A firm with $\sum_i\kappa_i>0$ is a provider of credit. The working capital of~\eqref{eq:ccc}
is $W = (\text{DSO}+\text{DIO}-\text{DPO})\times r/365$.

\section{Size and the direction of credit}

\subsection{The implication tested}

\begin{quote}
\textbf{Registered implication}: the smaller the size class, the larger
$\mathrm{DSO}-\mathrm{DPO}$; that is, smaller firms are the providers of credit.
\end{quote}

The ground was the claim of Section~\ref{sec:layers} that the credit layer is an image of
bargaining power: small firms without bargaining power should be the ones kept waiting for
payment.

\subsection{Result}

\begin{center}
\begin{tabular}{@{}lrrrr@{}}
\toprule
Size (capital) & DSO & DIO & DPO & DSO\textminus DPO \\
\midrule
1bn yen and over & 71 & 20 & 44 & \textbf{27} \\
100m -- 1bn & 61 & 20 & 47 & 15 \\
50m -- 100m & 47 & 20 & 37 & 10 \\
20m -- 50m & 49 & 15 & 39 & 10 \\
10m -- 20m & 53 & 18 & 35 & 18 \\
Under 10m & 35 & 12 & 21 & \textbf{14} \\
\bottomrule
\end{tabular}
\captionof{table}{DSO, DIO, DPO and DSO\textminus DPO by capital class.}
\label{tab:dso-dpo-by-size}
\end{center}
\begin{center}\footnotesize All industries excluding finance and insurance, FY2024, in days\end{center}

\begin{figure}[htbp]
\centering
\insertfig{ccc-trend}
\caption{DSO and DPO by capital class (all industries, FY2024). Both are short for the small classes.}
\label{fig:size-dso-dpo}
\end{figure}

\textbf{The implication is rejected; the sign is the opposite.} Firms of 1bn yen and over sit
at 27 days and those under 10m at 14, so large firms are the providers of credit. The ordering
is consistent over the ten years 2015--2024, and the large-firm figure has in fact widened
from 19 to 27 days.

\subsection{The error in the reasoning}

The error can be located. \textbf{It was overlooked that DSO and DPO both move.}

Small firms are kept waiting on receivables (DSO works against them). But at the same time
they are not trusted by suppliers either and cannot stretch payables. As
Figure~\ref{fig:size-dso-dpo} shows, the DPO of the under-10m class is 21 days, less than half
the 44 of firms of 1bn and over.

A lack of bargaining power acts on both sides: collection is slow but payment is fast. On
balance, \textbf{weak bargaining power does not appear in the net credit position}.

\begin{proposition}[Extending credit requires capital]
\label{prop:credit-needs-equity}
A net provider of credit, that is a party with $\sum_i\kappa_i>0$, satisfies $E > A$.
\end{proposition}

\begin{proof}
By the fundamental inequality of Proposition~\ref{prop:fundamental}, $W\leq E-A$. Since
$W=\sum_i\kappa_i+\mathrm{Inv}$ with $\mathrm{Inv}\geq 0$, $\sum_i\kappa_i>0$ implies $W>0$
and hence $E-A>0$.
\end{proof}

Under the $E\approx 0$ of Chapter~\ref{part:solo}, the requirement was the reverse,
$\kappa<0$ (Corollary~\ref{cor:zero-equity}). Tolerating collection as long as a DSO of 71
days is possible because there is capital behind it.

\subsection{The sign splits by industry}

\begin{center}
\begin{tabular}{@{}lrrr@{}}
\toprule
Industry & 1bn and over & Under 10m & Difference \\
\midrule
Information and communications & 64 & 15 & +49 \\
Construction & 67 & 19 & +48 \\
All industries & 27 & 14 & +13 \\
Manufacturing & 26 & 24 & +2 \\
Employment placement and dispatch & 24 & 25 & \textminus 1 \\
Retail & \textminus 16 & 6 & \textbf{\textminus 22} \\
\bottomrule
\end{tabular}
\captionof{table}{DSO\textminus DPO by industry, comparing capital of 1bn yen and over with under 10m. Only retail reverses sign.}
\label{tab:dso-dpo-by-industry}
\end{center}
\begin{center}\footnotesize DSO\textminus DPO, FY2024, in days\end{center}

\textbf{Only retail reverses sign.} The $-16$ days of large retailers is exactly the negative
credit position described in Part~\ref{part:theory}: cash collected immediately and purchases
paid in arrears. Here large firms are the recipients of credit and small retailers the
providers.

The $+49$ days of information and communications and construction shows large firms bearing
long advances in the position of contractor --- the territory of families 1-3 (one-off
arrears) and 5-3 (cost plus) of Part~\ref{part:catalog}.

The claim that ``larger firms are the providers of credit'' is therefore itself only an
average across industries. What decides is not size but \textbf{the $\Phi$ dominant in that
industry}. The lesson recorded in Part~\ref{part:results} --- that levels of $\kappa$ differ
structurally across industries so cross-sectional comparison is meaningless --- holds on the
size axis too.

\subsection{Implication for Chapter~\texorpdfstring{\ref{part:solo}}{VII}}

Chapter~\ref{part:solo} deduced that a solo business without assets requires $\kappa_C<0$.
The data show that requirement is not met: the under-10m class has DSO of 35 days and DPO of
21, a net 14 days as a provider.

The form differs from what was envisaged at registration, however. The small classes are not
extremely $\kappa>0$; rather \textbf{both DSO and DPO are short}. The transactions themselves
are small and short. Rather than lying outside the feasible region, they appear to sit at
different coordinates.

\section{Is $\CCC$ a leading indicator of the cycle?}
\label{sec:e-verification}

\subsection{The implication tested}

\begin{quote}
\textbf{Registered implication}: the rate of change of $\CCC$ \emph{leads or is synchronous
with} the rate of change of sales. In particular, a worsening receivables turnover
\emph{leads} a slowdown in sales.
\end{quote}

Correlations between the annual change $\Delta\CCC(t)$ and the sales growth rate $g(t+k)$
were computed by industry and aggregated over 45 industries, excluding aggregate and
short series.

\subsection{Result}

\begin{center}
\begin{tabular}{@{}lrrrl@{}}
\toprule
Lag & Mean correlation & Median & Share negative & Sign test \\
\midrule
Leading $k=-1$ & 0.036 & 0.058 & 47\% & $p=0.72$ \\
Synchronous $k=0$ & \textminus 0.159 & \textminus 0.228 & 69\% & $p=0.008$ \\
Lagging $k=+1$ & 0.072 & 0.118 & 31\% & --- \\
Lagging $k=+2$ & 0.085 & 0.035 & 44\% & --- \\
\bottomrule
\end{tabular}
\captionof{table}{Correlation between $\Delta\CCC(t)$ and sales growth $g(t+k)$ by lag, over 45 industries.}
\label{tab:ccc-sales-lag-corr}
\end{center}

\textbf{Synchrony is supported.} A worsening $\CCC$ and a slowdown in sales occur in the same
fiscal year: 31 of 45 industries show a negative correlation, with $p=0.008$ on a sign test.

\textbf{Leading is rejected outright.} 47\% negative is no information at all ($p=0.72$).
Change in $\CCC$ carries nothing about the following year's sales.

\subsection{The synchronous negative correlation may be a denominator effect}

The mechanism is likely the reverse of what was assumed. $\CCC = W/r$ has sales in the
denominator. If sales fall, then unless inventory and receivables fall at the same rate,
$\CCC$ worsens \emph{by definition}.

When the panel by medium for family 2 was treated in Part~\ref{part:results}, this was stated
explicitly as the alternative: issuance shrinks and the ratio rises mechanically.
Section~\ref{sec:denominator-effect} separates it here and shows that
\textbf{the denominator effect contributes about 17\%}.

\begin{center}
\begin{tabularx}{\textwidth}{@{}Xl@{}}
\toprule
Claim & Verdict \\
\midrule
Change in $\CCC$ is synchronous with change in sales & supported; the denominator effect contributes 17\% (Section~\ref{sec:denominator-effect}) \\
Change in $\CCC$ leads change in sales & \textbf{rejected} ($p=0.72$) \\
$\CCC$ is a leading indicator of the cycle & \textbf{rejected} \\
\bottomrule
\end{tabularx}
\captionof{table}{Verdicts on three claims about the leading property.}
\label{tab:predE-verdict}
\end{center}

\subsection{That the sign splits by industry is a finding}

Thirteen of the 45 industries have $|r|>0.4$ in the synchronous correlation.

\begin{center}
\begin{tabular}{@{}lrlr@{}}
\toprule
Strongly negative & $r$ & Strongly positive & $r$ \\
\midrule
Pulp, paper and paper products & \textminus 0.78 & Agriculture and forestry & +0.80 \\
Fabricated metal products & \textminus 0.67 & Other transport equipment & +0.69 \\
Information and communications & \textminus 0.58 & Accommodation and food services & +0.54 \\
Production machinery & \textminus 0.53 & Agriculture, forestry and fisheries & +0.52 \\
\bottomrule
\end{tabular}
\captionof{table}{Industries with $|r|>0.4$ in the synchronous correlation (four strongly negative, four strongly positive).}
\label{tab:sync-corr-extremes}
\end{center}

The strongly negative industries are all make-to-order with large inventories and
receivables, matching the conditions under which the denominator effect is strong.

What matters is that positive industries exist.
\textbf{A pure denominator effect would make every industry negative.} The positive
industries have a structure in which $\CCC$ grows as sales grow --- they build working capital
in order to grow --- which is the region where the $g^\star$ of~\eqref{eq:gstar} binds.

Change in $\CCC$ therefore contains components other than the denominator effect. They cannot
be separated in these data.

\section{An observation not anticipated: the long-run rise in \texorpdfstring{$\CCC$}{CCC}}

Separately from the tests, a fact not envisaged in advance was observed.

\begin{figure}[htbp]
\centering
\insertfig{lag-corr}
\caption{$\CCC$ and the net credit position, all industries excluding finance and insurance.}
\label{fig:ccc-trend}
\end{figure}

As Figure~\ref{fig:ccc-trend} shows, the $\CCC$ of all industries rose from 24.9 days in
FY2000 to 37.2 in FY2024, almost monotonically over 25 years. The net credit position also
doubled, from 9.8 to 18.5 days.
\textbf{The Japanese corporate sector as a whole is moving towards holding more working
capital.}

\begin{remark}[What the aggregate measures]
By Corollary~\ref{cor:kappa-nets-out}, credit positions closed within the corporate sector
cancel on aggregation. The $\CCC$ here is therefore not an indicator of inter-firm credit but
the sum of the position against the non-corporate sector and inventory, divided by sales.
Indeed, DSO $-$ DPO for all industries in Table~\ref{tab:dso-dpo-by-size} is 13 days, and
most of the gap to $\CCC=37.2$ days is inventory.
\textbf{The dominant component of the aggregate $\CCC$ is inventory.}
The same reason explains why negative $\CCC$ in individual firms does not appear in
aggregate.
\end{remark}

\subsection{The growth constraint has not worsened, however}

It is tempting to conclude from this that the growth constraint has worsened, since the
denominator of~\eqref{eq:gstar} has risen by half. \textbf{That inference is wrong.}

Over the same period the operating margin $m$ improved by 91\%, from 2.62\% to 5.01\%,
outrunning the worsening denominator. The ratio $m/\CCC$ consequently rose from 38.5\% to
49.2\% (Proposition~\ref{prop:gstar-flat}).

\begin{center}
\begin{tabular}{@{}lrrr@{}}
\toprule
 & FY2000 & FY2024 & Change \\
\midrule
$\CCC$ & 24.9 days & 37.2 days & +49\% \\
$m$ & 2.62\% & 5.01\% & +91\% \\
$m/\CCC$ & 38.5\% & 49.2\% & +28\% \\
\bottomrule
\end{tabular}
\captionof{table}{Change in $\CCC$, $m$ and $m/\CCC$ from FY2000 to FY2024.}
\label{tab:ccc-m-ratio-2000-2024}
\end{center}

What is used here is $m/\CCC$, not the $g^\star=(m+d-I/r)/\CCC$ of
Corollary~\ref{cor:gstar} (Remark~\ref{rem:gstar-flat-caveat}), because $d$ and $I$ were not
obtained far enough back.

\begin{remark}[Discipline in handling ratios]
$\CCC$ is the denominator of a ratio and says nothing on its own about the growth constraint.
Drawing a conclusion from a rise in $\CCC$ alone is looking at one half of a ratio.
\textbf{Concluding from one half of a ratio is the same kind of error as the retrospective
narrative that Section~\ref{sec:narrative} warns against.}
The difference is that the latter is an interpretation added after the fact, whereas this is
filling a shortfall in measurement with inference.
\end{remark}

The cause of the rise in $\CCC$ itself remains untested. Which of industrial composition,
regulation of payment terms, the decline of trade bills, or the level of interest rates
contributes is not separated here. The importance of the question is limited, however, since
$g^\star$ has not worsened.

\section{Levels by industry (for reference)}
\label{sec:ccc-extremes-by-industry}

As the lesson of Part~\ref{part:results} has it, cross-sectional comparison of levels cannot
be interpreted; the figures are recorded as a reference on the distribution of $\Phi$.

\begin{center}
\begin{tabular}{@{}lrlr@{}}
\toprule
Long $\CCC$ & Days & Short $\CCC$ & Days \\
\midrule
Leasing & 210 & Food services & 5 \\
Real estate & 89 & Water transport & 7 \\
Chemicals & 80 & Other transport & 9 \\
Textiles & 78 & Entertainment & 9 \\
General-purpose machinery & 77 & Motor vehicles and parts & 9 \\
Production machinery & 64 & Electricity & 10 \\
\bottomrule
\end{tabular}
\captionof{table}{Industries with the longest and shortest $\CCC$ (top six each).}
\label{tab:ccc-extremes-by-industry}
\end{center}
\begin{center}\footnotesize FY2024, excluding aggregate series\end{center}

The ordering agrees with the types of Part~\ref{part:catalog}. Leasing (family 3-1) and real
estate acquire an asset first and recover it over time; food services (family 1-1, immediate
exchange) sits at $\kappa\approx 0$. The structure of the differences by industry described
earlier appears directly.

\section{Decomposing the rise in \texorpdfstring{$\CCC$}{CCC}}
\label{sec:ccc-decomposition}

The rise in $\CCC$ is decomposed into its causes. The following implication was fixed in
advance.

\begin{quote}
\textbf{Registered implication}: the rise in $\CCC$ owes more to the composition effect than
to within-industry change, because industries with long $\CCC$ have gained weight.

\textbf{Alternative}: if the within-industry effect dominates, working capital has increased
in each industry irrespective of composition.
\end{quote}

\subsection{Shift-share decomposition}

By Corollary~\ref{cor:ccc-weighted}, the overall $\CCC$ is a sales-weighted mean of industry
figures, $\CCC=\sum_k w_k\CCC_k$. Its difference splits in two.
\begin{equation}
\Delta\CCC = \underbrace{\sum_k \bar{w}_k\,\Delta\CCC_k}_{\text{within-industry}}
 + \underbrace{\sum_k \Delta w_k\,\bar{\CCC}_k}_{\text{composition}}
\label{eq:shift-share}
\end{equation}
Taking $\bar{\cdot}$ to be the midpoint of the two periods,
\eqref{eq:shift-share} \textbf{holds exactly with no residual}. Substituting
$\bar w=(w_1+w_0)/2$ and $\bar\CCC=(\CCC_1+\CCC_0)/2$, the right-hand side equals
$\sum_k(w_{k1}\CCC_{k1}-w_{k0}\CCC_{k0})$. This is not an approximation introduced for
measurement but a consequence of $\CCC$ being a ratio
(Corollary~\ref{cor:ccc-weighted}).

\begin{center}
\begin{tabular}{@{}lrrr@{}}
\toprule
Period & $\Delta\CCC$ & Within-industry & Composition \\
\midrule
2000$\to$2024 (33 common industries) & +11.4 days & +9.5 days (82.9\%) & +1.9 days (17.1\%) \\
2010$\to$2024 (45 common industries) & +10.1 days & +9.9 days (98.1\%) & +0.2 days (1.9\%) \\
\bottomrule
\end{tabular}
\captionof{table}{Shift-share decomposition of $\Delta\CCC$ into within-industry and composition effects.}
\label{tab:shift-share-decomposition}
\end{center}

\textbf{The implication is rejected.} The within-industry effect accounts for 83\% to 98\%.
Since 2010, where the classification is stable, the composition effect is nearly zero. An
explanation by the shift towards services does not hold.

The reason the composition effect is small is also clear. Information and communications rose
in weight from 0.2\% to 5.8\%, pushing $\CCC$ up by $+3.05$ days, while electrical machinery
shrank from 6.3\% to 2.1\% ($-1.62$ days) and advertising from 5.6\% to 0.8\% ($-1.59$ days):
\textbf{the pushes up and down nearly cancel}.

\subsection{Components of the within-industry effect}

The within-industry effect splits into $\mathrm{DSO}$, $\mathrm{DIO}$ and $\mathrm{DPO}$.

\begin{center}
\begin{tabular}{@{}lrrrl@{}}
\toprule
Period & DSO & DIO & DPO & Main cause \\
\midrule
2000$\to$2024 & \textminus 0.93 & +3.49 & \textbf{+6.90} & shortening DPO \\
2010$\to$2024 & \textbf{+5.26} & +4.02 & +0.62 & lengthening DSO \\
\bottomrule
\end{tabular}
\captionof{table}{The within-industry effect split into DSO, DIO and DPO.}
\label{tab:within-industry-decomposition}
\end{center}
\begin{center}\footnotesize In days of contribution to $\CCC$\end{center}

Against Corollary~\ref{cor:ccc-change}, DSO and DPO correspond to $\tau_C$ and $\tau_S$ and
DIO to $\tau_{\mathrm{inv}}$. Hence:

\begin{center}
\begin{tabular}{@{}lrr@{}}
\toprule
Period & Change in contracts $\sum\Delta\tau_i$ & Change in operations $\Delta\tau_{\mathrm{inv}}$ \\
\midrule
2000$\to$2024 & +5.97 days & +3.49 days \\
2010$\to$2024 & +5.88 days & +4.02 days \\
\bottomrule
\end{tabular}
\captionof{table}{The split into change in contracts and change in operations, by Corollary~\ref{cor:ccc-change}.}
\label{tab:contract-vs-operation-change}
\end{center}

\textbf{In both periods, change in contracts accounts for about 60\%.} The rise in $\CCC$
originates in changes to contractual terms rather than in demand or inventory management.

\textbf{The main cause differs by period}, and the path of trade bills explains it.

\begin{center}
\begin{tabular}{@{}lrrr@{}}
\toprule
 & 2000 & 2010 & 2024 \\
\midrule
Notes receivable & 12.1 days & 6.1 days & 4.7 days \\
Notes payable & 15.3 days & 7.2 days & 5.1 days \\
\bottomrule
\end{tabular}
\captionof{table}{Notes receivable and payable, FY2000, FY2010 and FY2024.}
\label{tab:tegata-trend}
\end{center}

Bills fell to below a third during the 2000s. Notes payable lost 8.1 days, from 15.3 to 7.2,
and the 6.9-day shortening of DPO is almost entirely explained by this.
\textbf{Bills are a device for creating $\kappa_S<0$, and their disappearance pushed $\CCC$ up.}

In the 2010s the decline of bills had run its course and DSO lengthened by 5.26 days. The
breakdown is notes receivable $-1.37$ days and accounts receivable $+6.62$ days. Although bills
keep falling, accounts receivable grow by more than enough to offset them. By industry,
$\mathrm{DSO}$ rose in 28 of 45 industries, with the largest contributions from wholesale
(59$\to$70 days), retail (22$\to$28 days), information and communications (62$\to$78 days) and
construction (65$\to$71 days).

\begin{remark}[Not separable from a shift between accounts]
\label{rem:densai}
The growth of accounts receivable admits two readings: a genuine lengthening of the collection
period, and a shift between accounts accompanying the move to electronically recorded monetary
claims. Under the latter, the rise in $\CCC$ would be only apparent.

Separating them requires the balance of electronically recorded claims, which is not among the
survey items of the Financial Statements Statistics. Under the classification of
Section~\ref{sec:obstacle-taxonomy} this belongs to (iii) non-identification.
\end{remark}

\section[Separating the contemporaneous negative correlation]{Separating the contemporaneous negative correlation: the contribution of the denominator effect}
\label{sec:denominator-effect}

For the contemporaneous negative correlation observed in Section~\ref{sec:e-verification}, we
separate information about demand from the denominator effect inherent in the definition
$\CCC=W/r$.

In log differences, $\Delta\ln\CCC = \Delta\ln W - \Delta\ln r$. Under a pure denominator effect
$W$ would not move. Estimating the elasticity
\begin{equation}
\beta = \frac{\Delta \ln W}{\Delta \ln r}
\label{eq:wc-elasticity}
\end{equation}
therefore makes $1-\beta$ the contribution of the denominator effect.

\textbf{The theoretical value of $\beta$ is $1$.} By \eqref{eq:ccc}, $W=\CCC\cdot r$, so if
$\CCC$ is constant then $\Delta\ln W=\Delta\ln r$. Hence
\begin{equation}
\Delta\ln\CCC = (\beta-1)\,\Delta\ln r
\label{eq:beta-ccc}
\end{equation}
and the departure of $\beta$ from $1$ is exactly the sales elasticity of $\CCC$.
Equation~\eqref{eq:ccc-change} of Corollary~\ref{cor:ccc-change} decomposes this movement into
change in contracts and change in operations.

\begin{quote}
\textbf{Registered implication}: $\Delta\ln W$ and $\Delta\ln r$ are positively correlated; that
is, working capital falls when sales fall.
\end{quote}

\subsection{Result}

\begin{center}
\begin{tabularx}{\textwidth}{@{}Xr@{}}
\toprule
Pooled estimate (44 industries, 946 observations) & $\beta = 0.830$ \\
Coefficient of determination & $R^2 = 0.076$ \\
Median of industry-level $\beta$ & 0.625 \\
Standard deviation of industry-level $\beta$ & 1.161 \\
\midrule
Contribution of the denominator effect $1-\beta$ & 0.170 \\
\bottomrule
\end{tabularx}
\captionof{table}{Elasticity of $\Delta\ln W$ with respect to $\Delta\ln r$ (equation~\eqref{eq:wc-elasticity}).}
\label{tab:wc-elasticity-result}
\end{center}

\textbf{The implication is supported.} Working capital moves almost proportionally with sales,
and only about 17\% is attributable to the denominator effect. The contemporaneous negative
correlation of Section~\ref{sec:e-verification} carries, for the most part, information from the
demand side.

Substituting $\beta=0.830$ into \eqref{eq:beta-ccc} gives a sales elasticity of $\CCC$ of
$-0.170$: \textbf{$\CCC$ rises in phases where sales fall}. The amplification factor of
Remark~\ref{rem:ccc-amplifier} is then larger, so the degree to which a downturn in demand is
transmitted to cash strengthens in recessions. Calling $1-\beta$ the ``denominator effect'' is a
computational label, not a claim that $\CCC$ has not moved.

\subsection{Lag structure}
\label{sec:lag-test}

By Proposition~\ref{prop:kappa-tau}, $W(t)$ corresponds to the integral of $r$ over the past
$\tau$ periods. In industries with long $\CCC$ it should also depend on the previous year's
sales. The following was registered in advance.

\begin{quote}
\textbf{Registered implication (1)}: the coefficient $\beta_1$ on $\Delta\ln r(t-1)$ is larger the
longer that industry's $\CCC$.

\textbf{Registered implication (2)}: the total elasticity $\beta_0+\beta_1$ including the lag
exceeds the contemporaneous-only estimate of $0.830$, so the contribution of the denominator
effect is smaller than 17\%.
\end{quote}

The distributed-lag model
$\Delta\ln W(t)=\beta_0\Delta\ln r(t)+\beta_1\Delta\ln r(t-1)+\varepsilon$ was estimated on 886
observations across 43 industries.

\begin{center}
\begin{tabular}{@{}lr@{}}
\toprule
$\beta_0$ (contemporaneous) & 0.836 \\
$\beta_1$ (one-period lag) & \textminus 0.015 \\
Total elasticity & 0.821 \\
$R^2$ & 0.081 \\
\bottomrule
\end{tabular}
\captionof{table}{Estimates of the distributed-lag model.}
\label{tab:lag-model-result}
\end{center}

\textbf{Both implications are rejected.} $\beta_1$ is nearly zero, and its correlation with
$\CCC$ is $-0.101$ — the sign is also reversed. In industries with $\CCC<30$ days,
$\beta_1=+0.217$; in those with $\CCC\geq 60$ days, $-0.267$. The total elasticity falls slightly
from $0.830$ to $0.821$, leaving the contribution of the denominator effect essentially unchanged
at 17\% to 18\%.

\begin{remark}[Proposition~\ref{prop:kappa-tau} is not rejected]
\label{rem:lag-undetectable}
Failure to detect a lag structure does not mean equation~\eqref{eq:kappa-tau} is wrong. The level
of $\CCC$ is at most 210 days and mostly 30--60 days, which fits within a year at the granularity
of annual data. For $\beta_1$ to be significant, $\CCC$ would have to exceed one year, and no such
industry exists.

Quarterly data might allow detection, but the quarterly survey rests on the provisional closings
of sampled corporations and is therefore affected by closing adjustments.
\end{remark}

\begin{remark}[The confidence is limited]
$R^2=0.076$ is extremely low: 92\% of the variance of $\Delta\ln W$ is unexplained by
$\Delta\ln r$. The dispersion across industries is also large — $\beta<0.5$ in 14 industries,
$\beta>1.0$ in 13, and negative in some (retail $-0.43$, leasing $-0.01$).

Annual $\Delta\ln W$ mixes in year-end inventory adjustment, timing shifts of large transactions,
and changes in accounting policy. \textbf{One may say that the denominator effect is not dominant,
but this is insufficient to assert positively that the movement is information about demand.}
\end{remark}

\section{The effect of the scope of inventories on the conclusions}
\label{sec:inventory-scope}

The $\mathrm{DIO}$ of this chapter uses only ``finished goods or merchandise'' and excludes
work in process and raw materials. Since the Financial Statements Statistics carries
``inventories (end of period)'' as a separate item, the difference between the two measures the
size of the understatement. This section measures that size.

\textbf{This section is not a test of an implication.} It substitutes a different measurement
definition on the same data and checks how much the conclusions move.

\subsection{Two cross-sections}

``Inventories (end of period)'' was obtained along two cuts. \textbf{The two identify different
variation}, and must be read separately.

\begin{center}
\begin{tabularx}{\textwidth}{@{}lXr@{}}
\toprule
Cut & Coverage, and the variation identifying $\beta$ & Obs. \\
\midrule
Size cut & 10 industries $\times$ 11 size strata $\times$ FY2015--2024. $\beta$ is identified by variation \textbf{between size strata} & 1,100 \\
Industry cut & 62 industries $\times$ all sizes $\times$ FY2000--2024. $\beta$ is identified by variation \textbf{between industries and years} & 1,364 \\
\bottomrule
\end{tabularx}
\captionof{table}{The two cuts used to obtain inventories.}
\label{tab:inventory-cuts}
\end{center}

The estimate of this chapter (Section~\ref{sec:denominator-effect}) uses 946 observations across
45 industries, \textbf{the same design as the industry cut}. The figures that correct this chapter
are therefore those of the industry cut; the size cut answers a different question. Unless stated
otherwise, the figures below are from the industry cut.

\subsection{The level is understated roughly twofold}

The ratio of $\mathrm{DIO}$ has a median of 1.99, and the median $\CCC$ rises from 28.0 to 43.9
days. For all industries (excluding finance and insurance), FY2000 goes from 24.9 to 38.2 days
and FY2024 from 37.2 to 53.5 days.

The size cut gives a median ratio of 1.89 and a difference with median 4.1 days and mean 8.4 days
— comparable in magnitude.

\begin{center}
\begin{tabular}{@{}lrrrr@{}}
\toprule
Industry & DIO goods & DIO inv. & CCC goods & CCC inv. \\
\midrule
Manufacturing & 19.2 & 50.4 & 41.6 & \textbf{72.8} \\
Construction & 9.7 & 33.8 & 39.8 & \textbf{63.9} \\
All industries (excl.\ finance and insurance) & 18.7 & 35.0 & 37.2 & \textbf{53.5} \\
Information and communications & 7.5 & 14.6 & 58.1 & 65.2 \\
Other services & 8.0 & 12.2 & 34.4 & 38.6 \\
Wholesale & 24.2 & 27.6 & 32.9 & 36.3 \\
Food services & 3.4 & 6.1 & 4.9 & 7.7 \\
Accommodation & 6.5 & 8.9 & 21.8 & 24.2 \\
Retail & 24.4 & 25.4 & 23.3 & 24.2 \\
Employment placement and worker dispatch & 0.2 & 0.5 & 32.8 & 33.0 \\
\bottomrule
\end{tabular}
\captionof{table}{Difference in $\mathrm{DIO}$ and $\CCC$ by the scope of inventories (all sizes, FY2024, in days).}
\label{tab:inventory-scope}
\end{center}

The difference is of a different order across industries. $\CCC$ rises by more than 30\% in
manufacturing and construction, while retail and employment placement barely move.
\textbf{Divergence arises only in industries that hold work in process} — an obvious structure,
but large enough to change the ordering of the industry levels in
Section~\ref{sec:ccc-extremes-by-industry}.

\subsection{The effect on differences is small, but not negligible}

Comparing the two definitions for $\Delta\CCC$, the industry cut gives a correlation of 0.882 over
1,302 observations, with signs agreeing 86.8\% of the time. The size cut (990 observations) gives
a correlation of 0.976 and sign agreement of 91.0\%: \textbf{agreement is lower across
industries}. Because the weight of work in process differs by industry, including cross-industry
variation produces a gap.

\subsection{Correcting the definition worsens the elasticity slightly}

Equation~\eqref{eq:wc-elasticity} is estimated on a common sample. Since the observations with
$W>0$ differ by definition, \textbf{the comparison is restricted to observations with $W>0$ under
both definitions}.

\begin{center}
\begin{tabular}{@{}P{0.18\textwidth}P{0.30\textwidth}rrr@{}}
\toprule
Cut & Definition of $\mathrm{DIO}$ & $\beta$ & $R^2$ & Denom.\ effect $1-\beta$ \\
\midrule
Industry cut & finished goods or merchandise only & 0.840 & 0.073 & 0.160 \\
$n=955$ & inventories (total) & \textbf{0.828} & \textbf{0.166} & \textbf{0.172} \\
\midrule
Size cut & finished goods or merchandise only & 0.900 & 0.092 & 0.100 \\
$n=733$ & inventories (total) & 0.960 & 0.213 & 0.040 \\
\midrule
Estimate in the text & finished goods or merchandise only & 0.830 & --- & 0.170 \\
\bottomrule
\end{tabular}
\captionof{table}{Difference in the elasticity $\beta$ by the scope of inventories, on a common sample.}
\label{tab:inventory-elasticity}
\end{center}

In both cuts, taking inventories at their total more than doubles the coefficient of
determination. \textbf{Working capital is more nearly proportional to sales when work in process
and raw materials are included.}

\textbf{But the direction in which the denominator effect moves is opposite across the two cuts.}
In the size cut it falls from 0.100 to 0.040, while in the industry cut it rises slightly from
0.160 to 0.172.

\begin{remark}[Change the cut and the direction changes]
\label{rem:inventory-cut}
Correction shrinks the denominator effect in the size cut and enlarges it in the industry cut.
\textbf{The two are not error: they identify different variation}
(Table~\ref{tab:inventory-cuts}). Between size strata, larger firms hold more work in process, so
taking the total makes $W$ more nearly proportional to sales. Between industries, the weight of
work in process itself is dispersed as a characteristic of the industry, so taking the total does
not improve proportionality.

The estimate of this chapter is identified by variation between industries and years, so
\textbf{it is the industry cut that applies here}. A direction measured in one cut cannot be
carried into the other. \textbf{The range a sample covers is not the same as the range a
conclusion reaches.}
\end{remark}

\begin{remark}[Without a common sample the sign comes out reversed]
\label{rem:common-sample}
Since the observations with $W>0$ differ by definition, estimating without aligning the sample
gives $\beta$ of 0.894 and 0.787 in the size cut, yielding the opposite conclusion that the
denominator effect \emph{increases}. Restricted to the common sample the values are 0.900 and
0.960. \textbf{When measuring the effect of a definition, note that changing the definition
changes the composition of the sample.}
\end{remark}

\subsection{Verdict}

The conclusion of Section~\ref{sec:denominator-effect} is maintained, but
\textbf{it is merely maintained, not strengthened}.

\begin{center}
\begin{tabularx}{\textwidth}{@{}Xl@{}}
\toprule
Claim & Verdict \\
\midrule
The absolute level cannot be interpreted & Supported. The median $\CCC$ moves from 28.0 to 43.9 days \\
The effect on analyses of differences and correlations is small & Supported. But the correlation of $\Delta\CCC$ is 0.882, below the 0.976 of the size cut \\
The denominator effect is not dominant & Supported. But correction raises it slightly from 0.160 to 0.172 \\
\bottomrule
\end{tabularx}
\captionof{table}{Verdicts on three claims concerning the scope of inventories, by the industry cut.}
\label{tab:inventory-verdict}
\end{center}

All three are supported, but none strongly. That $1-\beta$ is 0.172 means 17\% of the increase in
working capital is not proportional to sales: not dominant, but not negligible either.

\section{An observation the simplified capacity constraint does not fit}
\label{sec:cap-observation}

Remark~\ref{rem:cap-simplification} introduced the simplification of treating $\mathrm{Cap}$ as a
constant. This section records a case in which it does not fit reality. It is a record of a
limitation, not a revision of the theory.

\subsection{The observation}

The price schedule of one major full-service fitness club chain was checked as of 2026. The
company offers a weekday-daytime-only plan for those aged 60 and over alongside its ordinary
plan. The two have identical structures — once a week (up to four visits a month), twice a week
(up to eight), and unlimited — and differ only in price.

\begin{center}
\begin{tabular}{@{}lrrr@{}}
\toprule
Facility category & Once/wk & Twice/wk & Unlimited \\
\midrule
I & 4.0\% & 10.3\% & 22.1\% \\
II & 4.8\% & 9.9\% & 21.5\% \\
III & 4.1\% & 9.6\% & 21.7\% \\
IV & 4.6\% & 10.0\% & 21.5\% \\
\bottomrule
\end{tabular}
\captionof{table}{Discount rate of the 60-and-over plan, by facility category and usage tier.}
\label{tab:senior-plan-discount}
\end{center}
\begin{center}\footnotesize Discount of the 60-and-over plan relative to the ordinary plan\end{center}

\textbf{In all four categories the discount increases monotonically with usage.} It stays around
4\% for once a week but reaches about 22\% on the unlimited plan. The price ratio from once a week
to unlimited is 1.67--1.74 on the ordinary plan and 1.36--1.43 on the 60-and-over plan.

\subsection{Inconsistency with the simplification}

Under equation~\eqref{eq:cap-simplified}, a large discount on the unlimited plan is hard to
explain, since it presses on the capacity constraint. Under the general form
\eqref{eq:cap-general}, however, it is explicable: restricting the plan to weekday daytime places
$\operatorname{supp}\delta_i$ within the idle part of $\mathrm{Cap}(t)$, so total volume can be
increased without worsening the constraint at the peak instant.

That is, \textbf{the operator manipulates the distribution over time, not the total}. The
simplification of this work cannot express this device.

\subsection{Why the theory is not revised from this observation}

\begin{remark}[No generalization from a case]
\label{rem:no-generalization}
The observation of this section rests on one company's price list.
Part~\ref{part:theory} is an abstraction intended to cover all $\Phi$, while
Part~\ref{part:results} is a concrete study confined to particular industries and firms.
\textbf{The concrete carries no claim of exhaustiveness.}

Rewriting equation~\eqref{eq:cap-simplified} from an $n=1$ observation would propagate a single
case across all 28 types. This is of the same form as the retrospective narrativization warned
against in Section~\ref{sec:narrative}: revising the structure of a population from one sample.

A revision of the theory is warranted when a deductive error is found, when the same limitation
appears in several independent domains, or when it affects the deductions of
Part~\ref{part:catalog}. None of these applies here. The observation is therefore recorded, and
no revision is made.
\end{remark}

\subsection{The simplified capacity constraint}

In connection with Section~\ref{sec:retiree}, the following implication was fixed in advance.

\begin{quote}
\textbf{Registered implication}: older cohorts choose pay-per-use over flat-rate contracts. This
follows from the preference predicted by socioemotional selectivity theory — that under a limited
time horizon one chooses the certain. A flat-rate contract is a bundle of options; pay-per-use is
certain.

\textbf{Alternative}: older cohorts have more disposable time and use the facility more often, so
the flat rate is more advantageous. Under this account the preference for flat rates
\emph{rises}.
\end{quote}

The signs are opposite, so the form is identifiable.

However, \textbf{no data on consumer choice could be obtained, and the implication is untested}.
What is needed is the distribution of contract-form choice by age, which operators hold but do not
publish.

Only the supply side's price design was obtained. The discount structure above rests on the
premise that older cohorts use the facility more, which is consistent with the alternative and
points against the registered implication. But \textbf{this is the operator's belief, not the
consumer's choice}. It requires the premise that the firm reads demand correctly, and so remains
indirect evidence.

\section{Limitations of this chapter}

\begin{enumerate}[label=(\arabic*)]
\item \textbf{The separation of the denominator effect is of low confidence.} The elasticity
estimate puts the contribution at about 17\%, but $R^2=0.076$, and most of the variation in
$\Delta\ln W$ is unexplained.

\item \textbf{The horizontal axis of the four quadrants is not measured.} No cross-industry
indicator exists for the ease of capturing surplus, so the classification remains one-dimensional,
on the demand side alone.

\item \textbf{The interpretation of the growth in accounts receivable is not settled.}
Section~\ref{sec:ccc-decomposition} excludes industrial composition and explains half by the
disappearance of bills, but whether the growth of accounts receivable in the 2010s is a genuine
lengthening or a shift between accounts cannot be separated (Remark~\ref{rem:densai}).

\item \textbf{The contract choice of older cohorts is untested.} As in
Section~\ref{sec:cap-observation}, no consumer-side choice data could be obtained. The supply
side's price design is consistent with the alternative but remains indirect evidence.

\item \textbf{Inventories are ``finished goods or merchandise'' only.} Work in process and raw
materials are excluded, so $\mathrm{DIO}$ is understated and the level of $\CCC$ comes out shorter
than it is. Section~\ref{sec:inventory-scope} measures the difference for ten industries and
confirms that the effect on analyses of differences and correlations is small, but
\textbf{correcting the 45-industry estimate would require obtaining that item for all
industries}.
\end{enumerate}

\section{On both implications being rejected}
\label{sec:two-rejections}

The direction of credit had the opposite sign; the leading property vanished. Both registered
implications failed.

Part~\ref{part:results} recorded the effect of pre-registration, and it operated more clearly in
this chapter. Had the implications not been fixed beforehand, the contemporaneous negative
correlation would very likely have been read as ``$\CCC$ works as a business-cycle indicator'' and
the credit position of large firms as ``large firms support their counterparties''. Both are
retrospective narratives.

The failure of the implications is not a rejection of the framework. The descriptions $\CCC$,
$\kappa$ and $\Phi$ function effectively; what failed were \textbf{conjectures about what these
quantities predict}. Effectiveness as a tool of description and effectiveness as a tool of
prediction are different things.
